\chapter{Conclusions and Future Work}
\label{ch:conclusions}

This chapter summarises the central empirical findings of the ablation study (Section~\ref{sec:conc:summary}), discusses the underlying methodological limitations (Section~\ref{sec:conc:limitations}), and outlines directions for subsequent research (Section~\ref{sec:conc:futurework}).

% ================================================================
\section{Summary of Findings}
\label{sec:conc:summary}
% ================================================================

This study posed three research questions, each operationalised as one objective of the controlled ablation in Chapter~\ref{ch:experiments}. This section answers each question in turn and then draws together the higher-level conclusions that follow from the three answers.

All three questions share a common premise established by the baseline (Section~\ref{sec:exp:baseline}): under MSE training, both architectures displace their predictions by roughly 50 to 55~minutes. Because the LSTM and PatchTST baselines exhibit almost the same delay despite their entirely different inductive biases, the delay cannot be attributed to the capacity of a particular architecture; it originates in the interaction between the pointwise loss and the strongly autocorrelated glucose signal (Section~\ref{sec:method:setup:cgm}). This is the failure mode that the three objectives set out to mitigate.

\paragraph{Research Question 1: do offset-aware loss functions reduce the delay and improve detection?}
The answer is largely negative. Both offset-aware losses reduce the mean lag only marginally relative to the baseline displacement of over fifty minutes, and for every combination of architecture and loss the lag-adjusted error RMSE$^*$ rises above the MSE baseline (Section~\ref{sec:exp:obj1}): the small temporal realignment is bought at the cost of trajectory shape rather than added to it. Soft-DTW, the loss from which most was expected, is the clearest case: it achieves the largest LSTM lag reduction yet the largest RMSE$^*$ increase, and the two architectures respond to it in qualitatively different ways. Tolerating or warping out timing errors in the loss therefore does not teach the model to anticipate change; the delay is not a defect of the optimisation criterion that a better-designed loss could remove.

\paragraph{Research Question 2: does multi-step supervision reduce the delay and improve detection?}
The answer is partial, and the way in which it is partial is itself informative. Multi-step supervision produces the most accurate trajectories of any configuration, with the lowest RMSE and MAE for the LSTM and the best values among the PatchTST variants (Section~\ref{sec:exp:obj2}). Yet the expectation that broader supervision would anchor the trajectory in time is not borne out: the mean lag falls by at most about one minute and event detection improves only slightly. The most accurate forecaster is thus not meaningfully less delayed than the baseline, the first direct evidence that trajectory quality and temporal alignment are dissociable.

\paragraph{Research Question 3: do direct classifiers outperform forecast-derived detection?}
The answer is affirmative for detection performance, with two qualifications that together form the most consequential result of the study. The direct classifiers achieve considerably higher recall and F2 than any forecast-derived configuration (Section~\ref{sec:exp:obj3:classifiers}): the detection gap is closed not by improving the forecast but by training directly on the detection objective. The tolerance sweep (Section~\ref{sec:exp:obj3:tausweep}) explains why the forecasting route is limited differently for the two architectures. Widening the acceptance window recovers a large share of PatchTST's detections, which are present in its output but displaced in time, whereas the LSTM's barely improve because it rarely produces a threshold crossing at all. Forecast-derived detection thus fails for a timing reason in PatchTST and a structural reason in the LSTM, which also explains why no single offset-aware loss helped both architectures in Objective~1.

The first qualification is an inversion of the architecture ranking: the LSTM, the weaker forecaster by RMSE$^*$, is the stronger classifier. The most plausible explanation lies in PatchTST's Reversible Instance Normalisation (Section~\ref{sec:bg:models:patchtst}), which removes the absolute glucose level that a threshold-referenced classifier depends on, whereas the LSTM operates on raw values and can learn proximity to the 70~mg/dL threshold directly. The second qualification comes from the lead-time analysis (Section~\ref{sec:exp:obj3:leadtime}): close to half of the LSTM classifier's true positives fire at zero lead time, registering an episode already under way rather than anticipating one, while the weaker PatchTST classifier issues a larger share of alarms that precede onset. Detection performance and anticipation are therefore distinct properties, and the F2 criterion, which rewards catching events but is indifferent to when the alarm is raised, captures only the first.

\paragraph{Synthesis.}
Taken together, the three answers support one overarching conclusion. The prediction delay cannot be optimised away at the level of the training loss or the output structure, since neither the offset-aware losses of Objective~1 nor the multi-step supervision of Objective~2 removes it; it is inherent to pointwise regression on a highly autocorrelated signal. Detection is improved not by mitigating the delay in the forecast but by replacing the forecasting objective with a detection objective (Objective~3); and the configuration that forecasts best (LSTM multi-step) is not the one that detects best (LSTM direct classifier), the ranking inverting between the two tasks. Glucose trajectory forecasting and hypoglycaemia event detection are therefore best treated as separate problems, each served by the model and objective suited to it. Even the strongest detection configuration, however, carries the two reservations established above: an F2 driven substantially by reactive rather than anticipatory alarms, and a precision low enough that false alarms remain frequent. These reservations motivate the directions set out below.

% ================================================================
\section{Methodological Limitations}
\label{sec:conc:limitations}
% ================================================================

Several structural constraints bound the conclusions of this study. First, and most fundamentally, all models operate on the univariate glucose signal alone, with no access to exogenous drivers such as insulin doses or carbohydrate intake. This deliberate choice isolates the temporal predictive capacity of the glucose signal itself, but it also qualifies the central finding: a rapid descent caused by a recent bolus or meal is, by construction, unpredictable from past glucose values, so part of the delay reflects missing causal information rather than a deficiency of the loss or the architecture. Whether the delay persists once such information is available is left open here, and is the first direction taken up below.

Second, the lead-time analysis shows that the LSTM classifier, despite achieving the highest F2, fires at zero lead time in nearly half of its true-positive detections. An alarm raised at the moment the threshold is crossed leaves no interval ahead of the event, so the F2 score overstates the configuration's value as an anticipatory detector; the F2 figures reported here should be read with this in mind.

Third, the event detection task is severely class-imbalanced, with only 5.7\% of labelled windows containing a positive event. This imbalance restricts the precision achievable across all configurations and means that even the highest-recall classifier produces a large volume of false alarms, a factor that would require dedicated mitigation before the classifier could be relied upon as a low-false-alarm detector.

Finally, the hyperparameters were tuned on a single designated patient rather than the full cohort, for computational reasons. Because all objectives share this one fixed configuration, the controlled-ablation comparisons are unaffected, since any sub-optimality applies equally to every variant; it bears only on absolute performance levels and the direct comparison between the two architectures, which should accordingly be read as indicative rather than definitive.

% ================================================================
\section{Future Work}
\label{sec:conc:futurework}
% ================================================================

Several directions follow from the limitations above. The most consequential is the integration of exogenous inputs, such as logged insulin doses and carbohydrate intake, which would lift the information ceiling that bounds the univariate approach and let the model anticipate excursions invisible to a glucose-only model. A lighter step, requiring no additional logging, would be to supply the time of day: meals and their accompanying insulin are among the strongest drivers of excursions~\citep{sun2026futureaware} and recur on an approximately daily schedule. Either extension would test whether the delay is genuinely irreducible or merely a consequence of withholding causal context.

A second theme concerns PatchTST's Reversible Instance Normalisation, which removes the absolute glucose level from the model's inputs (Section~\ref{sec:bg:models:patchtst}) and undercut two different configurations: it degraded amplitude accuracy when Soft-DTW aligned shapes in the scale-free space (Section~\ref{sec:exp:obj1}), and it deprived the PatchTST classifier of the absolute level that a 70~mg/dL threshold requires (Section~\ref{sec:exp:obj3:classifiers}). Two complementary steps follow: a scale-preserving Soft-DTW formulation, combining its elastic time alignment with a pointwise amplitude term or aligning before normalisation, would ensure the temporal realignment is no longer paid for in amplitude; and re-evaluating PatchTST with RevIN disabled or with a scale-preserving variant would test directly whether normalisation is responsible for its weak classification and for the architecture inversion observed in Objective~3.

A third direction is data-centric, targeting the class imbalance and the scarcity of informative segments. Rather than sampling uniformly over a signal that is mostly flat and event-free, training could concentrate on the most informative windows, those containing or immediately preceding a rapid descent toward the threshold. This could be combined with established imbalance remedies: a focal loss~\citep{lin2017focal} that down-weights easy negative examples, or oversampling and synthetic augmentation of the rare positive windows~\citep{chawla2002smote}. Together these aim to raise the precision side of the F2 score without sacrificing the recall the current classifiers already achieve.
